\section{Boh}

\[
    \sum_{i} \partial_{z_i} \left\langle \phi_1(z_1,\overline{z}_1) \cdots \phi_n(z_n,\overline{z}_n) \right\rangle = 0,
\]
\[
    \sum_{i} (z_i \partial_{z_i} + h_i) \left\langle \phi_1(z_1,\overline{z}_1) \cdots \phi_n(z_n,\overline{z}_n) \right\rangle = 0,
\]
\[
    \sum_{i} (z_i^2 \partial_{z_i} + 2 h_i z_i) \left\langle \phi_1(z_1,\overline{z}_1) \cdots \phi_n(z_n,\overline{z}_n) \right\rangle = 0,
\]
are the Ward identities for the global conformal transformations in two-dimensional conformal field theory. These identities impose constraints on the correlation functions of primary fields, which are fields that transform in a specific way under conformal transformations. The first identity corresponds to translations, the second to dilations, and the third to special conformal transformations. Solving these equations allows us to determine the form of correlation functions in two-dimensional CFTs, which is a crucial step in understanding the structure of these theories.

We can start from the two-point function of primary fields, which is constrained by conformal symmetry to take the form:
\[
    \begin{aligned}
        G(z_1,\overline{z}_1; z_2,\overline{z}_2)            & = \langle \phi_1(z_1,\overline{z}_1) \phi_2(z_2,\overline{z}_2) \rangle = \frac{C_{12}}{(z_1 - z_2)^{2h} (\overline{z}_1 - \overline{z}_2)^{2\overline{h}}}, \\
        G(\vert z_{12} \vert, \vert \overline{z}_{12} \vert) & = \frac{C}{\vert z_{12} \vert^{h_1 + h_2} \vert \overline{z}_{12} \vert^{\overline{h}_1 + \overline{h}_2}}, \quad z_{ij} = z_i - z_j,
    \end{aligned}
\]
where from the first line we see that the two-point function is non-zero only if the conformal weights of the two fields are equal, i.e., $h_1 = h_2$ and $\overline{h}_1 = \overline{h}_2$. This is a consequence of the Ward identities, which require that the correlation function be invariant under conformal transformations. If the weights were different, the correlation function would not be invariant under dilations or special conformal transformations, leading to a contradiction.

So the only two points correlation functions which can be different from zero are those in which the fields enter with the same weight; it may even be the same kind of operator, or a multiplet of operators with the same conformal weights. In this case, the two-point function is fixed up to a constant $C$. Indeed now, we can rewrite the constant as
\[
    C = \delta_{h_1, h_2} \delta_{\overline{h}_1, \overline{h}_2},
\]
whic leads us to the final form of the two-point function:
\[
    G(\vert z_{12} \vert, \vert \overline{z}_{12} \vert) = \frac{\delta_{h_1, h_2} \delta_{\overline{h}_1, \overline{h}_2}}{\vert z_{12} \vert^{h_1 + h_2} \vert \overline{z}_{12} \vert^{\overline{h}_1 + \overline{h}_2}} \simeq \frac{1}{(x_1-x_2)^{2\Delta}},
\]
where in the last step we have used the fact that the two-point function depends only on the distance between the two points, and we have defined $\Delta = h + \overline{h}$ as the scaling dimension of the field. This form of the two-point function is a direct consequence of conformal invariance and is a fundamental result in two-dimensional conformal field theory. It shows that the correlation functions of primary fields are highly constrained by symmetry, and it provides a starting point for understanding more complex correlation functions in these theories.

If we now insert the stress energy tensor $T(z)$ into the correlation function, we can derive the operator product expansion (OPE) between $T(z)$ and a primary field $\phi(w,\overline{w})$. The OPE takes the form:
\[
    \left\langle T(z) \phi(z_1,\overline{z}_1) \cdots \phi(z_n,\overline{z}_n) \right\rangle =\dots
\]
where we used the fact that the OPE of $T(z)$ with a primary field $\phi(w,\overline{w})$ is given by:
\[
    T(z) \phi(w,\overline{w}) \sim \frac{h}{(z-w)^2} \phi(w,\overline{w}) + \frac{1}{z-w} \partial_w \phi(w,\overline{w}) + \text{regular terms},
\]
where $h$ is the conformal weight of the primary field $\phi$. This OPE encodes how the stress energy tensor acts on the primary field, and it is a key ingredient in understanding the structure of conformal field theories. The first term in the OPE represents the scaling dimension of the primary field, while the second term represents the action of the Virasoro algebra on the primary field. Then we have a method to compute the correlation functions of primary fields by using the OPE with the stress energy tensor, which allows us to derive the Ward identities and ultimately determine the form of the correlation functions in two-dimensional conformal field theories.

We need correlation function to be invariant under the global conformal transformations, which are generated by the modes of the stress energy tensor. This leads to the Ward identities, which impose constraints on the correlation functions of primary fields.

For the three-point function, the Ward identities imply that it must take the form:
\[
    \left\langle \phi_1(z_1,\overline{z}_1) \phi_2(z_2,\overline{z}_2) \phi_3(z_3,\overline{z}_3) \right\rangle = \frac{C_{123}}{(z_{12})^{h_1 + h_2 - h_3} (z_{23})^{h_2 + h_3 - h_1}  (\overline{z}_{31})^{\overline{h}_1 + \overline{h}_3 - \overline{h}_2} (\overline{z} \setminus  \setminus )},
\]
where now the three-point function is fixed up to a constant $C_{123}$, which is known as the structure constant of the theory. The form of the three-point function is determined by conformal symmetry, and holds for all quasi-primary fields. The structure constant $C_{123}$ encodes the interactions between the three fields and is a fundamental quantity in conformal field theory: we need a method to find it.

Since solving a theory means solving or finding a way to compute all the correlation functions, we need to find a method to compute the structure constants. This is where the operator product expansion (OPE) comes into play. The OPE allows us to express the product of two fields as a sum over all possible fields in the theory, with coefficients that depend on the positions of the fields. By using the OPE, we can relate the three-point function to the two-point function and extract the structure constant $C_{123}$. It is not guaranteed that the OPE will still work in the same way as in higher dimensions, but in two-dimensional conformal field theories, the OPE is a powerful tool that allows us to compute correlation functions and determine the structure constants of the theory. This is a key step in solving two-dimensional conformal field theories and understanding their properties.






\section{Mode Expansion of the Stress Energy Tensor}

Starting from the conformal transformation in the \(z\) variable
\[
    w = z + \epsilon(z),
\]
we have the conserved charge associated with this transformation, which can be expressed as an integral of the stress energy tensor \(T(z)\) around a contour \(C\):
\[
    Q_{\epsilon} = \frac{1}{2\pi i} \oint_C \mathrm{d} z \, \epsilon(z) T(z).
\]

[...]

[...]


We can now compute the commutator of the modes \(L_n\) to find the Virasoro algebra. Starting from the definition of the modes:
\[
    L_n = \frac{1}{2\pi i} \oint_C \mathrm{d} z \, z^{n+1} T(z),
\]
we can compute the commutator \([L_m, L_n]\) starting from the product of two modes:
\[
    L_m, L_n = \dots
\]
which gives us this two integrals around the origin, with contraints on the radii of the contours, which then subtracted with the opposite order of integration, we can use singular part of the OPE of the stress energy tensor with itself (which we know from the previous calculation) to evaluate the integrals and find the commutation relations of the modes \(L_n\):
\[
    [L_m, L_n] = \oint_{0} \frac{\mathrm{d} w}{2\pi i} w^{m+1} \oint_{w} \frac{\mathrm{d} z}{2\pi i} z^{n+1} \left( \frac{c / 2}{(z-w)^4} + \frac{2 T(w)}{(z-w)^2} + \frac{\partial_w T(w)}{z-w} \right),
\]
where \(c\) is the central charge of the theory. To evaluate these integrals, we can use Cauchy formulae
\[
    \begin{aligned}
        f(z) = \oint_{z} \frac{\mathrm{d} w}{2 \pi i} \frac{f(w)}{w-z}, \\
        \frac{\mathrm{d}^n}{\mathrm{d} z^n} f(z) = n! \oint_{z} \frac{\mathrm{d} w}{2 \pi i} \frac{f(w)}{(w-z)^{n+1}},
    \end{aligned}
\]
to evaluate the integrals and find the commutation relations of the modes \(L_n\). We have to divide the calculation into three parts, corresponding to the three terms in the OPE of the stress energy tensor with itself.

The first term gives us a contribution proportional to the central charge \(c\),
\[
    \begin{aligned}
        I_1 & = \oint_{0} \frac{\mathrm{d} w}{2\pi i} w^{m+1} \oint_{w} \frac{\mathrm{d} z}{2\pi i} z^{n+1} \frac{c / 2}{(z-w)^4}                                          \\
            & = \dots                                                                                                                                                      \\
            & = \frac{c}{12}(m+1) m (m - 1) \oint \frac{\mathrm{d} w}{2\pi i} w^{m+n-1} = \frac{c}{12} (m^3 - m) \delta_{m+n,0} = \frac{c}{12} m (m^2 - 1) \delta_{n, -m},
    \end{aligned}
\]
where we have used the fact that the integral of \(w^{m+n-1}\) around the origin is zero unless \(m+n=0\), in which case it is equal to \(2\pi i\). The second term gives us a contribution proportional to the modes \(L_{m+n}\),
\[
    \begin{aligned}
        I_2 & = \oint_{0} \frac{\mathrm{d} w}{2\pi i} w^{m+1} \oint_{w} \frac{\mathrm{d} z}{2\pi i} z^{n+1} \frac{2 T(w)}{(z-w)^2} \\
            & = \frac{2}{2 \pi i} \oint_{0} \mathrm{d} w \, w^{m + n + 1} (m+1) T(w),
    \end{aligned}
\]
where if we consider the integral of \(w^{m+n+1} T(w)\) around the origin, we can use the definition of the modes \(L_k\) to express it in terms of the modes:
\[
    \oint_{0} \mathrm{d} w \, w^{m + n + 1} T(w) = 2\pi i L_{m+n},
\]
thus the second term gives us a contribution of
\[
    I_2 = 2 (m+1) L_{m+n}.
\]
For the third term, we have
\[
    \begin{aligned}
        I_3 & = \oint_{0} \frac{\mathrm{d} w}{2\pi i} w^{m+1} \oint_{w} \frac{\mathrm{d} z}{2\pi i} z^{n+1} \frac{\partial_w T(w)}{z-w} \\
            & = \oint_{0} \frac{\mathrm{d} w}{2\pi i} w^{m+1 n+1} \partial_w T(w)                                                       \\
            & = - \oint_{0} \frac{\mathrm{d} w}{2 \pi i} \partial_w (w^{m+n+2}) T(w) = - (m+n+2) L_{m+n},
    \end{aligned}
\]
where we have used integration by parts to evaluate the integral. Summing up the contributions from the three terms, we find the commutation relations of the modes \(L_n\):
\[
    [L_m, L_n] = (m-n) L_{m+n} + \frac{c}{12} (m^3 - m) \delta_{n, -m},
\]
which is the \textbf{Virasoro algebra}. This is the algebra of the modes generating the \(z\) part of the conformal transformations in two-dimensional conformal field theory, and it plays a central role in the structure of these theories. We have another virasoro algebra for the \(\overline{z}\) part of the conformal transformations, generated by the modes \(\overline{L}_n\) of the anti-holomorphic stress energy tensor \(\overline{T}(\overline{z})\), which satisfy the same commutation relations as the \(L_n\) modes:
\[
    [\overline{L}_m, \overline{L}_n] = (m-n) \overline{L}_{m+n} + \frac{c}{12} (m^3 - m) \delta_{n, -m}.
\]
Furthermore, since the holomorphic and anti-holomorphic parts of the stress energy tensor commute with each other, we also have
\[
    [L_m, \overline{L}_n] = 0.
\]
Thus the conformal group in two dimensions is generated by two copies of the Virasoro algebra, one for the holomorphic part and one for the anti-holomorphic part, and these two algebras commute with each other. This rich algebraic structure is a key feature of two-dimensional conformal field theories and allows for a deep understanding of their properties and classification.

The fields defined on the 2D worldsheet transforms under the action of the Virasoro algebra, and we can classify them according to their transformation properties. The first understanding of these transformations came from string theory, where the fields defined on the worldsheet of the string transform under the action of the Virasoro algebra, which is the algebra of the conformal transformations on the worldsheet.

This is an example of a Lie Algebra, but it is an infinite-dimensional Lie algebra, which is a key feature of two-dimensional conformal field theories. Its internal product is given by the commutator of the modes \(L_n\), satisfying the Jacobi identity. From these, mathematicians started studying anc classifying the infinite-dimensional Lie algebras, and the Virasoro algebra is one of the most important examples in this classification.

Any state in the Hilbert space and any field in the algebra of local fields in CFT must be classified in irreducible representations of the Virasoro algebra, which are called \textbf{Virasoro representations}. We have a huge hilbert space and operator algebra, but we can classify all the states and fields in terms of the representations of the Virasoro algebra, which is a powerful tool for understanding the structure of two-dimensional conformal field theories. Since the algebra is infinite dimensional too, we can have a finite or infinite number of irreducible representations (infinite dimensional), and it should diide in direct sums of irreducible representations.

The strategy wil be to use the representation theory of the Virasoro algebra to classify the fields and states in the theory, and then use this classification to compute correlation functions and other physical quantities. This is a key step in solving two-dimensional conformal field theories and understanding their properties.

We can see that Witt algebra is the algebra of the modes of the stress energy tensor without the central extension, i.e., without the central charge term in the commutation relations. The Virasoro algebra is a central extension of the Witt algebra, which means that it includes an additional term proportional to the central charge \(c\) in the commutation relations. The Witt algebra is a subalgebra of the Virasoro algebra, and it corresponds to the case where the central charge \(c\) is zero. In two-dimensional conformal field theories, the central charge \(c\) plays a crucial role in determining the properties of the theory, and it can take on different values depending on the specific model being studied.

The cenrak charge term, furthermore, disappears for \(m=-1,0,1\), so that if we restrict ourselves the Mobius transformations, which are generated by the modes \(L_{-1}, L_0, L_1\), we recover the Witt algebra without the central extension. This is consistent with the fact that the global conformal transformations in two dimensions form a finite-dimensional subgroup of the full infinite-dimensional conformal group, and this subgroup is generated by the modes \(L_{-1}, L_0, L_1\) of the Virasoro algebra (or Witt in this case).

    [... interpretations of the central charge, etc. ...]

